\section{Calibration findings}\label{sec:findings}

The previous section audits the environment. This one runs policies through it and reports what the resulting instruments read. Five findings follow: the baseline board under the corrected scoring kernel and the eligibility witness that answers the question the empty board raises (\cref{sec:f1,sec:witness}), regret against the closed-form market-making reference (\cref{sec:f2}), the generalization gap and cross-regime transfer matrix (\cref{sec:f3}), the failure-mode distribution (\cref{sec:f7}), and the replicated ecology probe (\cref{sec:f8}). Every follow-up unit sits directly after the finding it extends.

F1, F3, F7, F8 and the witness run on the canonical position-trading generator, whose tape the realism diagnostic of \cref{sec:f4} declines to certify on 23 of 24 seeded panels; F2 runs on the Avellaneda--Stoikov market-making model instead. Every number below is therefore a calibration of an instrument on tape the paper does not claim is market-realistic, and \cref{sec:limitations} carries that constraint forward into every claim built on it. The byte-identity guarantee covers the engine outputs behind F1 through F3 and the canonical linear path; the Python evidence reductions are seed-deterministic on one platform but sit outside the golden-hash path. F1 and F3 report 2{,}000-resample seed bootstrap intervals, F2 reports t-based 95\% intervals over committed per-episode vectors, and F7 and F8 are count-based, reporting episode counts and seed-level winner distributions rather than interval estimates. Every entry point, seed set and quantity is named where it is reported, and the commands are in \cref{app:commands}.

\subsection{F1: Pipeline validation under the corrected kernel}\label{sec:f1}

We roll the reference-policy field (flat, equal-weight long, momentum, minimum-variance, maximum-Sharpe, fractional-Kelly volatility targeting) over sixteen scenario seeds in each of the Calm, Hard and Extreme tiers via \texttt{run\_baselines}. The SharpeBench \texttt{score\_run} kernel scores each policy's pooled return series with the declared trial count equal to the field size, and each row carries a seed-paired bootstrap 95\% confidence interval on the deflated Sharpe (2{,}000 resamples).

This experiment validates the producer-scorer pipeline end to end; the scorer-side discovery it reproduces belongs to the companion paper. Under the pre-0.5.0 kernel this board was all zeros: every baseline scored a deflated Sharpe of 0.0000 on every tier, which reflected a unit error rather than strictness. The kernel applied its annualized 0.5 deflation prior per period without conversion, which set the bar near an annualized Sharpe of 18 on daily bars, a level nothing clears. Finding and fixing that unit bug is the companion SharpeBench paper's Finding 1 \citep{toca2026sharpebench}, corrected by converting once at 252 periods per year; what this paper claims is only that the producer-scorer separation surfaced it, because trajectories produced here could be re-scored under both kernels and the discrepancy attributed to the scorer alone.

\begin{table}[t]
\centering
\caption{Baseline leaderboard under the corrected kernel: deflated Sharpe (DSR) and pass\textsuperscript{$k$} rate per tier, 16 seeds each. No policy is rank-eligible on any tier: eligibility requires the per-run gate to pass on every seed (a pass\textsuperscript{$k$} rate of exactly 1.00), and the best rate is 0.75.}
\label{tab:f1}
\small
\setlength{\tabcolsep}{4.5pt}
\begin{tabular}{l rr rr rr}
\toprule
 & \multicolumn{2}{c}{Calm} & \multicolumn{2}{c}{Hard} & \multicolumn{2}{c}{Extreme} \\
\cmidrule(lr){2-3}\cmidrule(lr){4-5}\cmidrule(lr){6-7}
Policy & DSR & pass\textsuperscript{$k$} & DSR & pass\textsuperscript{$k$} & DSR & pass\textsuperscript{$k$} \\
\midrule
kelly\_vol\_target & 1.0000 & 0.75 & 0.0277 & 0.31 & 0.0029 & 0.12 \\
equal\_weight\_long & 1.0000 & 0.62 & 0.1114 & 0.31 & 0.0234 & 0.06 \\
min\_variance & 0.9849 & 0.44 & 0.0024 & 0.19 & 0.0243 & 0.06 \\
flat & 0.0007 & 0.00 & 0.0007 & 0.00 & 0.0007 & 0.00 \\
momentum & 0.0000 & 0.00 & 0.0000 & 0.00 & 0.0000 & 0.06 \\
max\_sharpe & 0.0000 & 0.06 & 0.0000 & 0.00 & 0.0000 & 0.00 \\
\bottomrule
\end{tabular}
\end{table}

\Cref{tab:f1} and \cref{fig:f1} show the corrected board. On Calm, drift is free deflated Sharpe: \texttt{kelly\_vol\_target} and \texttt{equal\_weight\_long} saturate the DSR at 1.0000, with bootstrap CIs $[0.9798, 1.0000]$ and $[0.8343, 1.0000]$. \texttt{min\_variance} reaches 0.9849 on a near-vacuous interval of $[0.0264, 1.0000]$. The deflated Sharpe is a probability bounded in $[0,1]$, so 1.0000 is its ceiling rather than an unbounded score, and the percentile bootstrap undercovers at that boundary: the saturated intervals are indicative, and no bias-corrected variant is applied. Yet nothing on the board is rank-eligible, because no policy passes the per-run gate on every seed; the best pass\textsuperscript{$k$} rate is 0.75. On Hard the best DSR collapses to 0.1114 and on Extreme to 0.0243, and the long baseline's pass\textsuperscript{$k$} rate degrades monotonically with stress (0.62, 0.31, 0.06). Each rate is 16 Bernoulli seeds, so near a rate of one half a single rate carries a nominal 95\% half-width of about $0.24$, and a difference between two of them about $0.35$: only the Calm-to-Extreme drop of $0.56$ exceeds that, so the monotone ordering describes these seeds rather than a resolved ranking. \texttt{momentum} and \texttt{max\_sharpe} lose money outright on every tier (mean per-bar returns down to $-1.6 \times 10^{-3}$ on Extreme).

Eligibility is a conjunction, and each leg catches what the other misses. Deflation alone would crown drift on Calm; pass\textsuperscript{$k$} alone would miss the overfit breadth the deflation discounts. The number a real agent has to beat is a positive deflated Sharpe that survives every evaluation seed, and no reference policy produces one. The board alone does not show that the eligibility set is non-empty for this task family. \Cref{sec:witness} answers that by construction, injecting an out-of-band oracle signal of controlled strength and locating the strength at which the gates first open; the eligibility set is non-empty on every tier, and the per-run gate that fails every baseline here is the leg that binds there too.

\begin{figure}[t]
\centering
\includegraphics[width=\linewidth]{f1-baselines}
\caption{Deflated Sharpe per baseline per tier with seed-paired bootstrap CIs. Calm saturates the long baselines; Hard and Extreme collapse them toward zero.}
\label{fig:f1}
\end{figure}

\subsubsection{Rank-eligibility witness: the eligibility set is non-empty}\label{sec:witness}

The baseline board above contains no rank-eligible policy, leaving open whether the conjunction can be satisfied at all. We answer by construction with an out-of-band oracle rather than an entrant. For each seed, the script replays the deterministic, policy-invariant tape one bar ahead and mixes standardized next-bar truth $z_t$ with a seeded noise path $\varepsilon_t$:
\[
\mathrm{signal}_t(s)=s z_t+\sqrt{1-s^2}\,\varepsilon_t,
\qquad 0\le s\le1.
\]
Within one noise path the same $\varepsilon_t$ is reused at every $s$ (common random numbers), so changing signal strength does not silently change the experiment. A sign-following policy trades every bar; a deadband policy trades only when $|\mathrm{signal}|>1$ and otherwise holds. Both are scored through the unchanged SharpeBench conjunction over 16 seeds. The per-seed gate uses the probabilistic Sharpe ratio (PSR), the probability under the kernel's stated sampling model that Sharpe exceeds its reference threshold. The script records the full coarse eligibility set, the sampled strengths at which the kernel conjunction accepts. It refines an \emph{observed crossing bracket} by bisection only when that coarse set is monotone; otherwise it reports the observed set without calling any crossing a threshold. A crossing located under one noise path is a functional of that draw, so we repeat the whole procedure (coarse sweep, monotonicity check, bisection) under five independent noise paths per cell: the first is the path serialized as the primary witness, the other four redraw $\varepsilon_t$. Each bisection bracket is numerical resolution on the stated grid (half-width $0.0016$ in $s$); the spread of the crossing across noise paths is a separate, statistical quantity, and it is what \cref{tab:witness} reports. All 50 attained coarse series (10 cells, five paths each) are monotone.

\begin{table}[t]
\centering
\caption{Observed eligibility crossing in signal-truth correlation $s$: mean [min, max] of the bisection-bracket midpoint over five independent noise paths per cell. The min to max range is variation across noise realizations; the bisection bracket around each individual crossing has half-width $0.0016$ and is not shown. Crossings are found by bisection after a monotone coarse-grid check on every path, and are not global optima or theorem-level thresholds. ``none (0/5)'' means that no sampled $s\le1$ is eligible on any of the five paths. At every one of the 50 attained crossings, the last failing gate is the requirement that every seed's marginal PSR reach $0.90$.}
\label{tab:witness}
\small
\setlength{\tabcolsep}{3.5pt}
\begin{tabular}{llccc}
\toprule
Variant & Band & Calm & Hard & Extreme \\
\midrule
\texttt{sign\_follow} & gap-band witness & none (0/5) & $0.730\ [0.702,0.752]$ & $0.343\ [0.302,0.389]$ \\
\texttt{sign\_follow} & F1 table & none (0/5) & $0.783\ [0.764,0.820]$ & $0.353\ [0.317,0.445]$ \\
\texttt{deadband\_hold} & gap-band witness & $0.848\ [0.777,0.930]$ & $0.500\ [0.408,0.623]$ & $0.337\ [0.283,0.383]$ \\
\texttt{deadband\_hold} & F1 table & $0.690\ [0.620,0.758]$ & $0.621\ [0.483,0.755]$ & $0.342\ [0.283,0.398]$ \\
\bottomrule
\end{tabular}
\end{table}

The result is an existence proof and a calibration. The deadband construction is eligible on every tier on every noise path, so the observed eligibility set is non-empty. The every-seed gate is last to open at all 50 observed crossings; pooled DSR and bootstrap gates open earlier. The crossing itself is not a constant of the tier: across noise paths it moves by $0.05$ to $0.27$ in $s$ (a factor of 15 to 85 above the bisection resolution), widest for the deadband policy on Hard, so any single-path threshold is one draw from that range and the table should be read at the precision of the range rather than of the bracket. Differences between the two bands are differences between two 16-seed scenario bands as well as between noise draws, and with five paths per cell we do not test them. The sign follower never becomes eligible on Calm, on any of the five paths, because its every-bar turnover costs more than the low-volatility signal is worth; at $s=1$ the noise term vanishes, so that endpoint is noise-invariant by construction, and the failure at every $s<1$ is what the replication adds. The environment therefore imposes no activity requirement, but its cost model does require a Calm policy to ration turnover (\cref{fig:witness}).

\begin{figure}[t]
\centering
\includegraphics[width=0.95\linewidth]{witness}
\caption{Witness sweep on the gap-band witness subset under the primary noise path. Filled markers are eligible sampled points; dashed lines mark the primary-path bisection crossing, drawn only after monotonicity is verified. Shaded bands (both panels) and the horizontal bars (right panels) give the min to max range and the mean of the crossing over five independent noise paths. The right panels show how many of the 16 seed-level PSR tests pass.}
\label{fig:witness}
\end{figure}

\subsection{F2: Regret against the closed-form reference policy}\label{sec:f2}

On the Avellaneda--Stoikov market-making task, \texttt{mm\_regret} runs each candidate quoting policy and the closed-form reference policy (\cref{sec:tasks}) on the same sixteen seeded episodes ($\sigma = 2.0$, $\gamma = 0.1$, $\kappa = 1.5$, 200 steps) and reports the mean reward gap to the reference. Candidates are the reference itself and fixed-spread quoters over a grid of half-spreads. The evidence commits the per-episode regret vector per grid point; intervals are t-based 95\% CIs over the sixteen paired episodes.

The reference's regret is 0.0 by construction, which confirms the plumbing. The fixed-spread curve is U-shaped (\cref{fig:f2}): regret is 84.6 (95\% CI $[68.2, 101.1]$) at a half-spread of 0.05 (quoting too tight, filled constantly, run over by inventory), falls to a minimum of 0.037 at 0.5, and rises again to 36.6 ($[33.0, 40.2]$) at 2.0 and 58.8 ($[55.3, 62.3]$) at 4.0 (quoting too wide, earning nothing). The minimum is the point the intervals do not resolve: the 0.5 point's CI is $[-8.7, 8.8]$ and the 1.0 point's is $[-5.4, 12.4]$. Both intervals include zero, so sixteen episodes do not resolve either point from the analytical reference; this failure to reject zero is not evidence of equivalence. At the observed per-episode dispersion, sixteen paired episodes detect a mean regret of about 12 reward units at 80\% power and a two-sided 5\% level, roughly a third of the 36.6 measured at a half-spread of 2.0, so differences smaller than that are outside this design's reach. The tight and wide extremes are separated from zero by wide margins. The instrument resolves the shape of the curve, not the sign of its minimum.

The point of shipping the reference is that agents on this task are scored on regret rather than raw PnL. A market-making leaderboard ranked on PnL rewards whichever policy drew the friendliest flow; regret against a fixed analytical reference on the same episodes removes that luck axis.

\begin{figure}[t]
\centering
\includegraphics[width=\linewidth]{f2-regret}
\caption{Mean regret versus the analytical Avellaneda--Stoikov reference policy over 16 paired episodes, with t-based 95\% CIs. Fixed-spread quoting is U-shaped in the half-spread; the reference defines zero.}
\label{fig:f2}
\end{figure}

\subsection{F3: Generalization gap and cross-regime transfer}\label{sec:f3}

\texttt{generalization\_gap} scores the default reference policy (equal-weight long) on disjoint train and held-out seed bands (16 seeds each, separated by a 10{,}000-seed gap) within each tier. \texttt{cross\_regime\_transfer} then holds sixteen seeds fixed and varies the regime, producing the full three-by-three transfer matrix whose diagonal is zero by construction, because identical regimes reuse byte-identical environments. The policy is fixed and parameter-free, so nothing is trained: ``source regime'' below means the regime the policy is nominally selected on, and for such a policy the matrix reduces to pairwise differences of three per-tier deflated Sharpes. What F3 measures is the instrument, how much score a regime switch moves for a fixed reference, and not transfer learning. These are the 16-seed subsets \cref{sec:protocol-bands} exempts fixed policies from replacing with the full canonical band. The evidence commits per-seed return series per cell; CIs are seed-resampled percentile bootstrap ($B = 2000$), seed-paired across regimes for the matrix cells.

\begin{table}[t]
\centering
\caption{Left: within-tier generalization gap (train minus test deflated Sharpe) for the fixed equal-weight reference. Right: cross-regime transfer gap (in-distribution minus zero-shot out-of-distribution DSR), rows source regime, columns scored regime; the diagonal is zero by construction. Cells are read only through their seed-paired bootstrap intervals, and the Hard$\leftrightarrow$Extreme cells ($\pm 0.096$), whose interval straddles zero, are not interpreted.}
\label{tab:f3}
\small
\begin{tabular}{lrrr}
\toprule
Tier & Train DSR & Test DSR & Gap \\
\midrule
Calm & 1.0000 & 0.9988 & $+0.0012$ \\
Hard & 0.1231 & 0.4618 & $-0.3388$ \\
Extreme & 0.0269 & 0.1280 & $-0.1012$ \\
\bottomrule
\end{tabular}
\hspace{1.5em}
\begin{tabular}{lrrr}
\toprule
 & Calm & Hard & Extreme \\
\midrule
Calm & 0.000 & $+0.877$ & $+0.973$ \\
Hard & $-0.877$ & 0.000 & $+0.096$ \\
Extreme & $-0.973$ & $-0.096$ & 0.000 \\
\bottomrule
\end{tabular}
\end{table}

\Cref{tab:f3} reports both. The within-tier gaps are the expected control: the reference policy has no fitted parameters, so it cannot overfit its train band, and the instrument correctly reads near zero on Calm ($+0.0012$) and noise-signed values on the stressed tiers (the negative Hard and Extreme gaps mean the held-out band happened to be kinder; with no learning, the gap is pure seed noise, and how much band luck a 16-seed evaluation carries is read from the interval widths rather than from any single realized gap). The per-band bootstrap intervals say the same thing at interval width: Hard's train-band DSR of 0.1231 carries a CI of $[0.000, 0.944]$ and its test band $[0.037, 0.948]$ (the bands are disjoint seed sets, so these are unpaired), overlapping almost entirely, which is what a pure-noise gap looks like. Hard's train-band DSR of 0.1231 here and \cref{tab:f1}'s 0.1114 for the same policy on the same sixteen seeds differ only in the declared trial count the deflation prior consumes, six declared policies in F1 against no declared search in the gap metric; rescoring the committed F3 return series at six trials reproduces 0.1114 exactly.

The transfer matrix (\cref{fig:f3}) is where regime structure appears, and every cell is classified by one rule: whether its seed-paired bootstrap interval excludes zero. The Calm-involved point estimates are the largest: source regime Calm scored zero-shot on Hard costs $0.877$ of deflated Sharpe, and Calm to Extreme costs $0.973$ of a score bounded in $[0,1]$. The seed-paired bootstrap intervals temper the two claims differently: Calm to Extreme excludes zero ($[0.635, 0.999]$), while Calm to Hard's interval is wide ($[0.054, 0.999]$), positive throughout but reaching close to zero, so at 16 seeds only the Calm-to-Extreme collapse is resolved. The Hard-to-Extreme cell ($0.096$, CI $[-0.218, 0.911]$) has an interval straddling zero, so it is not interpreted: at 16 seeds the instrument cannot distinguish it from band luck. A bounded, heavily deflated score bootstrapped over 16 seeds is a coarse instrument, which is itself a calibration result, and \cref{sec:protocol-seeds} turns it into a prescription. The matrix is antisymmetric because the policy is fixed, and its diagonal is zero by construction. Its off-diagonal magnitude is the quantity a within-tier gap is structurally blind to: a policy scored only on calm seeds can carry a saturated DSR into a leaderboard while being worth 0.03 under stress. Both metrics are cheap, and the protocol requires reporting both.

\begin{figure}[t]
\centering
\includegraphics[width=0.85\linewidth]{f3-transfer-matrix}
\caption{Cross-regime transfer matrix for the reference policy: out-of-distribution deflated Sharpe with the seed band held fixed.}
\label{fig:f3}
\end{figure}

\subsection{F7: Failure-mode distributions}\label{sec:f7}

\texttt{classify\_episode\_failure} classifies every rollout of the eight-policy field (the six reference policies plus the two behavioral counterparties), across all three tiers and sixteen seeds (384 episodes), into the worst-case-first taxonomy: cascade-wiped, bankrupt, stopped-out, mandate breach (structural, drawdown or inventory), or clean, with a per-scenario sampled mandate and a 50\% drawdown stop in force.

\begin{table}[t]
\centering
\caption{Failure-mode counts per tier, 128 episodes each (8 policies $\times$ 16 seeds). No episode reaches bankruptcy or a cascade wipe; failure is dominated by mandate breaches, and drawdown breaches triple from Calm to Extreme.}
\label{tab:f7}
\small
\setlength{\tabcolsep}{5pt}
\begin{tabular}{lrrrrrrr}
\toprule
Tier & Clean & Stopped out & Mandate struct. & Mandate DD & Mandate inv. & Bankrupt & Clean rate \\
\midrule
Calm & 72 & 0 & 36 & 11 & 9 & 0 & 0.56 \\
Hard & 68 & 0 & 36 & 14 & 10 & 0 & 0.53 \\
Extreme & 50 & 5 & 31 & 32 & 10 & 0 & 0.39 \\
\bottomrule
\end{tabular}
\end{table}

Mandates are sampled per scenario and assigned to fixed policies that cannot observe or adapt to them, so a structural breach is guaranteed whenever the draw pairs, say, a shorting policy with a long-only mandate. The flat structural-breach counts across tiers (36, 36, 31) are the signature of that assignment mechanism, not a behavioral finding about agents; a mandate-aware field could invert the distribution entirely, and measuring one is the interesting follow-up.

\Cref{tab:f7} and \cref{fig:f7} show the shape of failure as stress rises. The clean rate falls from 0.56 on Calm to 0.39 on Extreme. Each rate is 128 episodes, so a difference between two tiers carries a nominal binomial 95\% half-width of about $0.12$; the episodes sharing a seed or a policy are not independent, so that width is a floor. The Calm-to-Hard step (0.56 versus 0.53) sits inside it and is not interpreted; the drop to Extreme is the only step that clears it. Structural mandate breaches are roughly flat across tiers, as the mechanism above dictates: a shorting policy under a sampled long-only mandate breaks the rule on every tier. The stress-sensitive modes are the drawdown family: mandate drawdown breaches triple from 11 on Calm to 32 on Extreme, and hard stop-outs at the 50\% line appear only on Extreme (5 episodes, four of them \texttt{max\_sharpe}). No policy in this field ever reaches bankruptcy or a cascade wipe; the taxonomy's worst states are reachable but not by these baselines at this stop level. Per policy, \texttt{flat} is clean 48/48 by construction, while \texttt{overconfident} accounts for the most Calm failures (7 of its 16 episodes end in a drawdown breach) and \texttt{max\_sharpe} is the worst on Extreme (3 clean of 16). On this mandate-blind reference field, roughly half of the failures are process failures that a return-ranked board does not record.

\begin{figure}[t]
\centering
\includegraphics[width=\linewidth]{f7-failures}
\caption{Failure-mode distribution per tier per policy over 384 classified episodes.}
\label{fig:f7}
\end{figure}

\subsection{F8: Ecology outcomes under shocks, replicated over seeds}\label{sec:f8}

\texttt{run\_ecology} seeds the replicator with the eight baseline species over a shared-book market via \texttt{market\_payoffs} (12 generations, field size 8, innovation every 4 generations breeding parameter-perturbed variants of the leader) and runs trajectories under a steady all-Calm control schedule and a \texttt{regime\_shocks} schedule rotating Calm (generations 0 to 3), Hard (4 to 7) and Extreme (8 to 11). We replicate the probe over eight replicator seeds per schedule; \cref{tab:f8} reports the cross-seed winner distribution, and the seed-0 trajectories (\cref{fig:f8}) remain the illustrative detail. The two schedules chain episode seeds differently (the shock schedule strides seeds across generations), so control and shocked runs at the same replicator seed are matched in configuration, not in price paths.

The replicator flow rule sits in the market-ecology tradition: strategies as species whose profitability decays as their share of the flow grows \citep{farmer2002market}, and populations whose composition, rather than any single agent's optimization, drives aggregate dynamics \citep{lux1999scaling}. The implementation abstracts away much of what those models study: capital flows react instantly (no fund-flow lags), there are no leverage or margin constraints beyond the strategies' own sizing, and there is no entry beyond the single leader-perturbation innovator. Its outputs are trajectories of this replicator map, not claims about market ecology in general.

\begin{table}[t]
\centering
\caption{Replicated ecology outcomes over 8 replicator seeds per schedule: distribution of the final dominant lineage (bred variants credited to their founding species). The root winner differs between the schedules on 5 of 8 seeds, and on 3 of the 6 seeds where both schedules resolve a dominant lineage; the seed-0 replacement of the volatility targeter by the unlevered long book occurs on 1 of 8 seeds.}
\label{tab:f8}
\footnotesize
\setlength{\tabcolsep}{4pt}
\begin{tabular}{lrr}
\toprule
Final dominant lineage & Control (all Calm) & Shocked (Calm/Hard/Extreme) \\
\midrule
kelly\_vol\_target & 4 & 7 \\
equal\_weight\_long & 2 & 1 \\
flat (no resolved dominant) & 2 & 0 \\
\bottomrule
\end{tabular}
\end{table}

The replication does not support the single-seed narrative. At seed 0 the story was clean: \texttt{kelly\_vol\_target} sweeps the control board (peak share 1.00, final 0.81), collapses under the shock schedule once the regime turns Hard, and \texttt{equal\_weight\_long} inherits dominance (final 0.67). Across eight seeds that story is the exception. Under the shock schedule a \texttt{kelly\_vol\_target} lineage ends dominant on 7 of 8 seeds (final shares near 0.90); the seed-0 handoff to the long book happens once. The control side is itself seed-dependent: \texttt{kelly\_vol\_target} wins 4 seeds, \texttt{equal\_weight\_long} wins 2, and on 2 seeds no species resolves dominance at all (all eleven end persistent at small shares, with a \texttt{flat} variant holding the largest share at 0.10). The root-level winner differs between the two schedules on 5 of 8 seeds, but on 2 of those the control side never resolves a dominant lineage at all; among the 6 seeds where both schedules resolve one, the winner differs on 3. At eight seeds the reordering is neither established nor excluded, and it has no consistent direction. Extinction pressure is more stable than identity: shocked runs produce a resolved dominant on 8/8 seeds with 4 to 7 extinctions per run (at seed 0, both behavioral counterparties are gone or collapsed on both schedules, \texttt{overconfident} extinct at generation 0).

One replicator trajectory is one draw from a heavy-tailed outcome distribution, and single-run ecology narratives do not survive replication. What the probe reliably produces at this scale is the outcome distribution itself: which lineages can win, how often shocks reorder the board, and how much of the population goes extinct, per committed seed set. Sharper claims need more seeds, varied shock orderings and path-matched schedule pairs, all of which are mechanical extensions of the committed script.

\begin{figure}[tbp]
\centering
\begin{subfigure}[t]{0.92\linewidth}
\includegraphics[width=\linewidth]{f8-ecology-control}
\caption{Control: steady Calm.}
\end{subfigure}

\vspace{0.8em}

\begin{subfigure}[t]{0.92\linewidth}
\includegraphics[width=\linewidth]{f8-ecology-shocked}
\caption{Shocked: Calm, then Hard, then Extreme.}
\end{subfigure}
\caption{Seed-0 population-share trajectories over 12 replicator generations, the illustrative single run behind \cref{tab:f8}'s distribution. At this seed the regime rotation replaces the calm-regime winner with the unlevered long book; across the eight-seed replication that replacement occurs once.}
\label{fig:f8}
\end{figure}
